% Auxiliary fragment. Requires amsmath/amssymb and a standard article context.
\section{The ancestor-sensitive cost of removing the Shift halt}
This auxiliary does not meet the 194-state target. It separates a rejected
196-state shortcut from a correctly implemented 197-state source variation.
The parent is the MP196 prototype in the pinned incoming packet
\texttt{pa\_sub194\_research.zip}. Its successful modus ponens halts
immediately on consequent code one; otherwise it returns its consequent
and leaves the formula stack unchanged. Its source and arithmetic
calibration are inherited dependencies.

Write $p(a,b)=(a+b)(a+b+1)/2+b$, with stack encoding
$[a;w]=p(a,w)$. A command store $(t,w,d,h)$ is \emph{one-free} if $t$,
$h$ and all the canonical entries obtained by successively unpairing $w$
are different from one. The tail strictly decreases for nonzero stacks,
so the assertion is primitive recursive. The zero stack has further heads
zero, since $p(0,0)=0$.

The initialized store is one-free. An AX command pushes only its saved
one-free top and the generated implication or universal heads, which are
congruent to zero or two modulo four; guarded suppression gives zero.
An unsuccessful MP clears the visible fields, while a successful MP with
consequent one halts before a completed command boundary. Other successful
MP steps return a one-free consequent. A Shift at a one-free boundary has
$t+h\ne1$ and only rearranges existing formulas. Induction over completed
commands proves the invariant without assuming consistency of PA.

It is tempting to remove the delayed Shift test from
\[
 \mathtt{node21}=(\mathtt{node7},\mathtt{node20}),\qquad
 \mathtt{node7}=(\mathtt{node6},\mathtt{halt})
\]
by replacing the first expression by
$(\mathtt{node6},\mathtt{node20})$. This is incorrect.
The mismatch branch of \texttt{node6} contains an ancestor-sensitive
\texttt{breakTo 1}. With the wrapper present it exits \texttt{node7}
and continues into \texttt{node20}; with the wrapper removed it exits
\texttt{node21}, skipping the stack pop. At the one-free store
$t=h=0$, $w=p(2,0)=3$, the shortcut keeps $w=3,h=0$ instead of
producing $w=0,h=2$. The physical checker also fails its 175th attempted
contract, at source site 130 on the zero outcome of logical decrement 8.
Thus neither a source proof nor an implementation proof is available
for this 196-state shortcut.

For a correct invariant-based variation, retain the wrapper and replace
only its Shift occurrence of the halt leaf:
\[
 \mathtt{node21}=((\mathtt{node6},\mathtt{dec}\,5),\mathtt{node20}).
\]
The original \texttt{node7} remains in the immediate-MP test. The changed
leaf can be reached only if the old Shift comparison has succeeded,
namely $t+h=1$. It is therefore unreachable from any initialized
nonhalting command boundary. The old and new source executions agree
step for step on all reachable source configurations; their halting
behavior is equal. This is an invariant argument, not a claim that the
modified Shift has the usual semantics on arbitrary stores containing
one.

The literal implementation adds one reader to MP196:
\[
\begin{array}{c|cc}
 &0&1\\\hline
402&0R77&1R45
\end{array}
\qquad
(53,0):0R68\longmapsto0R402.
\]
The complete table consequently has 197 working controls and 394 rows.
Its 3,663 source sites give 4,769 primitive outcomes, all verified by
the inherited symbolic physical checker. The exact initial complete
frame is reached after 58,154 transitions, with native control 84,
head $-2$, logical registers $(2,0,0,0,0,0,0,0,0)$ and 26 one-valued
spare blocks. The generic arithmetic-box rows are unchanged. These
finite symbolic contracts, combined with the explicit register loop
proofs, supply the physical simulation; a new proof-assistant endpoint
is not claimed.

The no-one invariant therefore supplies a legitimate logical
simplification but does not save a control in this compilation. The
halt leaf already has zero working-control cost, while preserving its
ancestor scope and distinguishing its two occurrences costs an extra
reader. This is a precise obstruction to the displayed deletion, not
a lower bound for other source arrangements or Turing machines.
